\chapter{Sample Risk Assessment Templates}
\label{chap:Sample Risk Assessment Templates}

The templates in this appendix are provided as pro-forma working documents that construction organisations may adapt to their own projects, risk profiles, and management systems. Each template represents a minimum framework; organisations should add columns, rows, and fields as required by the specific activity and project context. All templates should be completed by a competent person with knowledge of the specific activity, the site conditions, and the applicable legal and guidance requirements. Completed templates should be retained for a minimum of three years (or longer where the activity involves substances hazardous to health, asbestos, or other long-latency health risks).

%---------------------------- SECTION ----------------------------
\section{Generic Site Risk Assessment Form}

This generic risk assessment form may be used for any construction site activity. The assessor should complete all fields; the Hazard and Who Affected columns should be specific to the activity and the site, not generic. Inherent Risk is the risk rating before controls; Residual Risk is the rating after controls are applied. Use a 5$\times$5 risk matrix (Likelihood 1--5, Severity 1--5; Risk = L$\times$S; Low 1--4, Medium 5--9, High 10--14, Very High 15--25).

\begin{table}[H]
\centering
\small
\begin{tabularx}{\textwidth}{>{\raggedright\arraybackslash}p{2.0cm}>{\raggedright\arraybackslash}p{2.2cm}>{\raggedright\arraybackslash}p{1.4cm}>{\raggedright\arraybackslash}X>{\raggedright\arraybackslash}p{1.4cm}>{\raggedright\arraybackslash}p{1.6cm}>{\raggedright\arraybackslash}p{1.5cm}}
\toprule
\textbf{Hazard} & \textbf{Who Affected} & \textbf{Inherent Risk L$\times$S} & \textbf{Control Measures} & \textbf{Residual Risk L$\times$S} & \textbf{Review Date} & \textbf{Assessor} \\
\midrule
Trench collapse & Groundworkers, banksman, adjacent operatives & 4$\times$5 = 20 & Hydraulic trench box installed; daily competent person inspection; no unsupported entry $>$1.2\,m; exclusion zone for plant. & 2$\times$5 = 10 & Monthly / after rainfall & J. Smith CSCS \\
\midrule
Buried service strike & Groundworkers, plant operators & 3$\times$5 = 15 & HSG47 sequence: utility records; CAT \& Genny survey; safe digging zone; hand-dig within 500\,mm; vacuum excavation at confirmed positions. & 1$\times$5 = 5 & Before each excavation & J. Smith CSCS \\
\midrule
Manual handling (pipe laying) & Groundworkers & 3$\times$3 = 9 & Mechanical assist for pipes $>$25\,kg; team lift for pipes 15--25\,kg; training; pre-task briefing. & 2$\times$3 = 6 & Monthly & J. Smith CSCS \\
\midrule
\textit{(Add rows as required for each identified hazard)} & & & & & & \\
\bottomrule
\end{tabularx}
\caption*{Generic Site Risk Assessment --- Project: \underline{\hspace{4cm}} Date: \underline{\hspace{2cm}} Assessed by: \underline{\hspace{3cm}}}
\end{table}

%---------------------------- SECTION ----------------------------
\section{Hot Work Permit to Work}

A Hot Work Permit to Work must be completed, authorised, and in place before any hot work activity commences. Hot work includes: welding (MMA, MIG, TIG); oxy-acetylene cutting; grinding that produces sparks; disc cutting; use of blow torches; and any other process that introduces a source of ignition into the workplace. The permit must be issued daily --- a multi-day hot work permit is not acceptable. The issuing authority must physically inspect the work location before issuing the permit.

\begin{table}[H]
\centering
\small
\begin{tabularx}{\textwidth}{>{\raggedright\arraybackslash}p{3.5cm}X}
\toprule
\multicolumn{2}{c}{\textbf{HOT WORK PERMIT TO WORK}} \\
\midrule
\textbf{Permit Number} & \underline{\hspace{6cm}} \\
\textbf{Project / Site} & \underline{\hspace{6cm}} \\
\textbf{Location of hot work} & \underline{\hspace{6cm}} (Building, floor, grid reference) \\
\textbf{Description of work} & \underline{\hspace{6cm}} \\
\textbf{Type of hot work} & $\square$ Welding \quad $\square$ Cutting \quad $\square$ Grinding \quad $\square$ Blow torch \quad $\square$ Other: \underline{\hspace{2cm}} \\
\textbf{Operative(s) names} & \underline{\hspace{6cm}} \\
\textbf{Permit valid from} & \underline{\hspace{2.5cm}} to \underline{\hspace{2.5cm}} (same day only) \\
\midrule
\multicolumn{2}{l}{\textbf{Pre-work fire precautions checklist (issuing authority to complete)}} \\
\midrule
Combustible materials removed or protected within 10\,m radius & $\square$ Yes \quad $\square$ N/A \\
Walls, partitions, and floors checked for voids and combustible linings & $\square$ Yes \quad $\square$ N/A \\
Extinguishers in place (type and number): \underline{\hspace{3cm}} & $\square$ Yes \\
Fire watch operative designated: \underline{\hspace{3cm}} & $\square$ Yes \\
Hot work area clear of flammable liquids, gases, and dusts & $\square$ Yes \\
Adjacent areas including floor below checked and clear & $\square$ Yes \\
Area of work ventilated or LEV in place & $\square$ Yes \quad $\square$ N/A \\
Passive fire protection (cavity barriers, intumescent seals) not compromised & $\square$ Yes \quad $\square$ N/A \\
\midrule
\textbf{Authorisation} & \\
Issuing Authority (Name / Signature) & \underline{\hspace{5cm}} Date: \underline{\hspace{2cm}} \\
Operative acknowledgement (Signature) & \underline{\hspace{5cm}} \\
\midrule
\multicolumn{2}{l}{\textbf{Post-work closure (to be completed by fire watch operative 60 minutes after completion)}} \\
Area inspected and free of smouldering material & $\square$ Yes \\
No heat detected in adjacent voids or surfaces & $\square$ Yes \\
Permit closed (Signature / Time) & \underline{\hspace{5cm}} Time: \underline{\hspace{2cm}} \\
\bottomrule
\end{tabularx}
\end{table}

%---------------------------- SECTION ----------------------------
\section{Confined Space Permit to Work}

A Confined Space Permit to Work is required before any entry into a space identified as a confined space under the Confined Spaces Regulations 1997. The permit system must be supported by a site-specific confined space rescue plan, pre-tested atmospheric monitoring equipment, and a trained rescue team in attendance throughout entry. The permit must be signed by both the issuing authority and the entrant; a separate copy must be displayed at the entry point and retained on file.

\begin{table}[H]
\centering
\small
\begin{tabularx}{\textwidth}{>{\raggedright\arraybackslash}p{3.5cm}X}
\toprule
\multicolumn{2}{c}{\textbf{CONFINED SPACE PERMIT TO WORK}} \\
\midrule
\textbf{Permit Number} & \underline{\hspace{6cm}} \\
\textbf{Location / Space description} & \underline{\hspace{6cm}} \\
\textbf{Nature of work} & \underline{\hspace{6cm}} \\
\textbf{Entrant(s)} & \underline{\hspace{6cm}} \\
\textbf{Standby person} & \underline{\hspace{6cm}} \\
\textbf{Rescue team members} & \underline{\hspace{6cm}} \\
\textbf{Permit valid} & From: \underline{\hspace{2cm}} To: \underline{\hspace{2cm}} (max 8-hour shift) \\
\midrule
\multicolumn{2}{l}{\textbf{Atmospheric testing record (prior to entry and at intervals during work)}} \\
\midrule
\textbf{Parameter} & \textbf{Pre-entry} \quad \textbf{30\,min} \quad \textbf{60\,min} \quad \textbf{Continuous monitor?} \\
Oxygen \% (acceptable: 19.5\%--23\%) & \underline{\quad\quad\quad\quad} \quad \underline{\quad\quad\quad\quad} \quad \underline{\quad\quad\quad\quad} \quad $\square$ Yes \\
Flammable gas \%LEL (acceptable: $<$10\%) & \underline{\quad\quad\quad\quad} \quad \underline{\quad\quad\quad\quad} \quad \underline{\quad\quad\quad\quad} \quad $\square$ Yes \\
CO ppm (acceptable: $<$20\,ppm) & \underline{\quad\quad\quad\quad} \quad \underline{\quad\quad\quad\quad} \quad \underline{\quad\quad\quad\quad} \quad $\square$ Yes \\
H\textsubscript{2}S ppm (acceptable: $<$1\,ppm) & \underline{\quad\quad\quad\quad} \quad \underline{\quad\quad\quad\quad} \quad \underline{\quad\quad\quad\quad} \quad $\square$ Yes \\
Gas monitor serial no. / calibration date & \underline{\hspace{5cm}} \\
\midrule
\multicolumn{2}{l}{\textbf{Emergency rescue plan (attach separate rescue plan document)}} \\
Rescue plan document reference & \underline{\hspace{5cm}} \\
Rescue equipment in place at entry point & $\square$ Tripod and winch \quad $\square$ Harness(es) \quad $\square$ BA \quad $\square$ Resuscitation kit \\
Emergency services pre-notified & $\square$ Yes \quad Location given: \underline{\hspace{2.5cm}} \\
\midrule
\multicolumn{2}{l}{\textbf{Entry log}} \\
\textbf{Name} & \textbf{Time in} \quad \textbf{Time out} \quad \textbf{Signature} \\
\underline{\hspace{3cm}} & \underline{\hspace{1.5cm}} \quad \underline{\hspace{1.5cm}} \quad \underline{\hspace{3cm}} \\
\underline{\hspace{3cm}} & \underline{\hspace{1.5cm}} \quad \underline{\hspace{1.5cm}} \quad \underline{\hspace{3cm}} \\
\midrule
\textbf{Authorisation (Issuing Authority)} & Name: \underline{\hspace{3cm}} Sig: \underline{\hspace{2cm}} Date: \underline{\hspace{1.5cm}} \\
\textbf{Permit closure} & Sig: \underline{\hspace{3cm}} Time closed: \underline{\hspace{2cm}} \\
\bottomrule
\end{tabularx}
\end{table}

%---------------------------- SECTION ----------------------------
\section{Electrical Isolation Permit to Work}

The Electrical Isolation Permit is required before any work is carried out on or near a de-energised electrical circuit or system. It ensures that the circuit cannot be inadvertently energised while work is in progress. The permit must be completed by a suitably qualified and competent person (typically an authorised person under the company's electrical safety rules). Lock-out Tag-out (LOTO) procedures must be applied where equipment is being isolated.

\begin{table}[H]
\centering
\small
\begin{tabularx}{\textwidth}{>{\raggedright\arraybackslash}p{3.5cm}X}
\toprule
\multicolumn{2}{c}{\textbf{ELECTRICAL ISOLATION PERMIT TO WORK}} \\
\midrule
\textbf{Permit Number} & \underline{\hspace{6cm}} \\
\textbf{Circuit / Equipment description} & \underline{\hspace{6cm}} \\
\textbf{Location} & \underline{\hspace{6cm}} \\
\textbf{Work to be carried out} & \underline{\hspace{6cm}} \\
\textbf{Authorised Person} & Name: \underline{\hspace{3cm}} Qualification: \underline{\hspace{2.5cm}} \\
\textbf{Working Person(s)} & \underline{\hspace{6cm}} \\
\midrule
\multicolumn{2}{l}{\textbf{Isolation and verification sequence}} \\
\midrule
1. Circuit / equipment identified on drawing or label & Drawing ref: \underline{\hspace{3cm}} $\square$ Confirmed \\
2. Isolation point identified (switch, MCB, fuse, disconnect) & Location: \underline{\hspace{3cm}} \\
3. Isolation method & $\square$ Switched off \quad $\square$ MCB off \quad $\square$ Fuse removed \quad $\square$ Disconnected \\
4. Lock-off applied & $\square$ Yes \quad Lock number: \underline{\hspace{2cm}} \\
5. Danger tag applied & $\square$ Yes \quad Tag text: \underline{\hspace{3cm}} \\
6. Test for dead (using approved voltage indicator) & Instrument: \underline{\hspace{2cm}} Cal date: \underline{\hspace{1.5cm}} \\
   & Phase L1: \underline{\quad} V \quad Phase L2: \underline{\quad} V \quad Phase L3: \underline{\quad} V \quad N-E: \underline{\quad} V \\
   & $\square$ Confirmed dead \\
7. Proving unit test (before and after) & $\square$ Instrument proved live before test \quad $\square$ Proved live after test \\
\midrule
\textbf{Permit issued} & Date: \underline{\hspace{2cm}} Time: \underline{\hspace{1.5cm}} Sig (AP): \underline{\hspace{3cm}} \\
\textbf{Permit accepted} & Sig (WP): \underline{\hspace{3cm}} \\
\midrule
\multicolumn{2}{l}{\textbf{Reinstatement}} \\
Work complete, tools and materials clear & $\square$ Confirmed \\
Isolation removed and circuit re-energised by AP & Date: \underline{\hspace{2cm}} Time: \underline{\hspace{1.5cm}} Sig (AP): \underline{\hspace{3cm}} \\
Permit cancelled & $\square$ Yes \\
\bottomrule
\end{tabularx}
\end{table}

%---------------------------- SECTION ----------------------------
\section{RAMS Cover Sheet}

The RAMS (Risk Assessment and Method Statement) Cover Sheet is the summary document that accompanies the detailed risk assessment and method statement for a specific activity or work package. It provides a single-page reference for the principal contractor's RAMS review process, distribution, and revision history. The cover sheet must be completed before the RAMS is submitted for approval; the RAMS must not commence use until the ``Approved'' box is signed.

\begin{table}[H]
\centering
\small
\begin{tabularx}{\textwidth}{>{\raggedright\arraybackslash}p{3.5cm}X}
\toprule
\multicolumn{2}{c}{\textbf{RISK ASSESSMENT AND METHOD STATEMENT --- COVER SHEET}} \\
\midrule
\textbf{Document title} & \underline{\hspace{8cm}} \\
\textbf{Document reference} & \underline{\hspace{4cm}} Version: \underline{\hspace{1.5cm}} \\
\textbf{Project / Contract} & \underline{\hspace{8cm}} \\
\textbf{Issuing organisation} & \underline{\hspace{8cm}} \\
\textbf{Prepared by} & Name: \underline{\hspace{3cm}} Role: \underline{\hspace{2cm}} Date: \underline{\hspace{1.5cm}} \\
\textbf{Reviewed by} & Name: \underline{\hspace{3cm}} Role: \underline{\hspace{2cm}} Date: \underline{\hspace{1.5cm}} \\
\midrule
\textbf{Scope of works} & (Brief description of activities covered by this RAMS) \newline \underline{\hspace{8cm}} \newline \underline{\hspace{8cm}} \\
\midrule
\multicolumn{2}{l}{\textbf{Hazard summary (principal hazards identified in this RAMS)}} \\
1. & \underline{\hspace{8cm}} \\
2. & \underline{\hspace{8cm}} \\
3. & \underline{\hspace{8cm}} \\
4. & \underline{\hspace{8cm}} \\
(See full risk assessment for complete hazard list) & \\
\midrule
\textbf{PPE required} & $\square$ Hard hat \quad $\square$ Hi-viz \quad $\square$ Safety boots \quad $\square$ Gloves \quad $\square$ Goggles \quad $\square$ RPE \quad $\square$ Harness \quad Other: \underline{\hspace{1.5cm}} \\
\midrule
\multicolumn{2}{l}{\textbf{Distribution}} \\
\textbf{Recipient} & \textbf{Role} \quad \textbf{Date issued} \quad \textbf{Signature of receipt} \\
\underline{\hspace{3cm}} & \underline{\hspace{2cm}} \quad \underline{\hspace{1.5cm}} \quad \underline{\hspace{3cm}} \\
\underline{\hspace{3cm}} & \underline{\hspace{2cm}} \quad \underline{\hspace{1.5cm}} \quad \underline{\hspace{3cm}} \\
\midrule
\multicolumn{2}{l}{\textbf{Revision history}} \\
\textbf{Version} & \textbf{Date} \quad \textbf{Description of revision} \quad \textbf{Approved by} \\
1.0 & \underline{\hspace{1.5cm}} \quad Initial issue \quad \underline{\hspace{3cm}} \\
\midrule
\textbf{RAMS Approval (Principal Contractor)} & Approved by: \underline{\hspace{3cm}} Role: \underline{\hspace{2cm}} \newline Signature: \underline{\hspace{3cm}} Date: \underline{\hspace{1.5cm}} \newline $\square$ Approved \quad $\square$ Approved with comments (see attached) \quad $\square$ Rejected (resubmit) \\
\bottomrule
\end{tabularx}
\end{table}

%---------------------------- SECTION ----------------------------
\section{Pre-Start Daily Inspection Checklist --- Excavations}

This checklist must be completed by a competent person before any operative enters an excavation at the start of each working day, after any rainfall or adverse weather, and after any event that may have affected the stability of the excavation. The competent person must have sufficient knowledge of ground behaviour, trench support, and service hazards to identify signs of incipient failure. Any defect identified must be rectified before operatives enter the excavation. All completed checklists must be retained on site for the duration of the project and for three years after project completion.

\begin{table}[H]
\centering
\small
\begin{tabularx}{\textwidth}{>{\raggedright\arraybackslash}p{0.6cm}>{\raggedright\arraybackslash}p{5.5cm}>{\raggedright\arraybackslash}p{1.2cm}>{\raggedright\arraybackslash}p{1.2cm}>{\raggedright\arraybackslash}p{1.2cm}X}
\toprule
\textbf{No.} & \textbf{Inspection Item} & \textbf{Satisfactory} & \textbf{Defect} & \textbf{N/A} & \textbf{Action / Notes} \\
\midrule
1 & Trench sides: no tension cracks visible at surface within 1.5$\times$ depth of trench & $\square$ & $\square$ & $\square$ & \\
2 & Trench sides: no progressive sloughing, bulging, or movement visible & $\square$ & $\square$ & $\square$ & \\
3 & Trench support: trench box / sheeting / piling in place, undamaged, and correctly positioned & $\square$ & $\square$ & $\square$ & \\
4 & Trench support: no displacement of waling frames, props, or packs & $\square$ & $\square$ & $\square$ & \\
5 & Water ingress: no significant water accumulation in trench; pump operational & $\square$ & $\square$ & $\square$ & \\
6 & Buried services: service markers in place; safe digging zone boundaries marked & $\square$ & $\square$ & $\square$ & \\
7 & Surcharge: no plant, vehicles, or spoil within surcharge exclusion zone (1.5$\times$ trench depth) & $\square$ & $\square$ & $\square$ & \\
8 & Access: safe means of entry and egress (ladder at no more than 9\,m intervals) & $\square$ & $\square$ & $\square$ & \\
9 & Exclusion zone: barriers erected at safe distance from trench edge & $\square$ & $\square$ & $\square$ & \\
10 & Gas detection: detector in use where made ground, near gas infrastructure & $\square$ & $\square$ & $\square$ & \\
11 & Adjacent structures: no sign of ground movement, cracking, or subsidence adjacent to trench & $\square$ & $\square$ & $\square$ & \\
12 & Inspection record: previous inspection defects resolved & $\square$ & $\square$ & $\square$ & \\
\midrule
\multicolumn{6}{l}{\textbf{Entry permitted:} $\square$ Yes \quad $\square$ No (defect(s) to be resolved first)} \\
\multicolumn{6}{l}{Inspector Name: \underline{\hspace{3cm}} Signature: \underline{\hspace{3cm}} Date: \underline{\hspace{1.5cm}} Time: \underline{\hspace{1.5cm}}} \\
\bottomrule
\end{tabularx}
\end{table}

%---------------------------- SECTION ----------------------------
\section{Pre-Start Daily Inspection Checklist --- Lifting Operations}

This checklist must be completed before the commencement of lifting operations on any working day. It covers cranes (mobile, tower, crawler), MEWPs, and hoists. The appointed person or lifting supervisor must verify that all items are in order before the first lift of the day. A separate lift plan must exist for each complex or non-routine lift; this checklist does not replace the lift plan.

\begin{table}[H]
\centering
\small
\begin{tabularx}{\textwidth}{>{\raggedright\arraybackslash}p{0.6cm}>{\raggedright\arraybackslash}p{5.5cm}>{\raggedright\arraybackslash}p{1.2cm}>{\raggedright\arraybackslash}p{1.2cm}>{\raggedright\arraybackslash}p{1.2cm}X}
\toprule
\textbf{No.} & \textbf{Inspection Item} & \textbf{Satisfactory} & \textbf{Defect} & \textbf{N/A} & \textbf{Action / Notes} \\
\midrule
1 & LOLER thorough examination certificate valid and in date & $\square$ & $\square$ & $\square$ & \\
2 & Operator's daily pre-use inspection completed and records available & $\square$ & $\square$ & $\square$ & \\
3 & Crane / MEWP / hoist safe working load (SWL/WLL) marked and legible & $\square$ & $\square$ & $\square$ & \\
4 & Lift plan for today's lifting operations reviewed and understood by operator and banksman & $\square$ & $\square$ & $\square$ & \\
5 & Ground conditions assessed; outrigger plates / mats in place for crane & $\square$ & $\square$ & $\square$ & \\
6 & Radius exclusion zone established and signed; unauthorised persons excluded & $\square$ & $\square$ & $\square$ & \\
7 & Banksman / slinger-signaller briefed and in position & $\square$ & $\square$ & $\square$ & \\
8 & Lifting accessories (slings, shackles, chains) inspected; colour-coded for current period & $\square$ & $\square$ & $\square$ & \\
9 & Load weight confirmed and within SWL/WLL for configuration & $\square$ & $\square$ & $\square$ & \\
10 & Overhead obstructions (power lines, structures, other cranes) identified and risk controlled & $\square$ & $\square$ & $\square$ & \\
11 & Wind speed within manufacturer's operating limit & $\square$ & $\square$ & $\square$ & \\
12 & Operator hold valid CPCS / IPAF ticket for the equipment & $\square$ & $\square$ & $\square$ & \\
\midrule
\multicolumn{6}{l}{\textbf{Lifting operations permitted:} $\square$ Yes \quad $\square$ No (defect(s) to be resolved first)} \\
\multicolumn{6}{l}{Lifting Supervisor Name: \underline{\hspace{3cm}} Signature: \underline{\hspace{3cm}} Date: \underline{\hspace{1.5cm}} Time: \underline{\hspace{1.5cm}}} \\
\bottomrule
\end{tabularx}
\end{table}

%---------------------------- SECTION ----------------------------
\section{Toolbox Talk Record}

The Toolbox Talk Record must be completed for every toolbox talk delivered on site. It provides the evidence that health and safety information was communicated to the workforce; the attendance list with signatures constitutes a legal record. Completed forms must be retained on site and archived for three years after the project end date.

\begin{table}[H]
\centering
\small
\begin{tabularx}{\textwidth}{>{\raggedright\arraybackslash}p{3.5cm}X}
\toprule
\multicolumn{2}{c}{\textbf{TOOLBOX TALK RECORD}} \\
\midrule
\textbf{Date} & \underline{\hspace{6cm}} \\
\textbf{Time} & \underline{\hspace{6cm}} \\
\textbf{Location / Site} & \underline{\hspace{6cm}} \\
\textbf{Topic of toolbox talk} & \underline{\hspace{6cm}} \\
\textbf{Presenter name and role} & \underline{\hspace{6cm}} \\
\textbf{Duration} & \underline{\hspace{3cm}} minutes \\
\midrule
\multicolumn{2}{l}{\textbf{Key messages delivered (summary):}} \\
\multicolumn{2}{X}{1. \underline{\hspace{10cm}}} \\
\multicolumn{2}{X}{2. \underline{\hspace{10cm}}} \\
\multicolumn{2}{X}{3. \underline{\hspace{10cm}}} \\
\midrule
\textbf{Questions raised} & \underline{\hspace{6cm}} \\
\textbf{Actions arising} & \underline{\hspace{6cm}} \\
\midrule
\multicolumn{2}{l}{\textbf{Attendees --- print name and sign to confirm attendance and understanding:}} \\
\multicolumn{2}{l}{\textbf{Name (print)} \hfill \textbf{Signature} \hfill \textbf{Employer / Trade}} \\
\multicolumn{2}{l}{\underline{\hspace{4cm}} \hfill \underline{\hspace{4cm}} \hfill \underline{\hspace{3cm}}} \\
\multicolumn{2}{l}{\underline{\hspace{4cm}} \hfill \underline{\hspace{4cm}} \hfill \underline{\hspace{3cm}}} \\
\multicolumn{2}{l}{\underline{\hspace{4cm}} \hfill \underline{\hspace{4cm}} \hfill \underline{\hspace{3cm}}} \\
\multicolumn{2}{l}{\underline{\hspace{4cm}} \hfill \underline{\hspace{4cm}} \hfill \underline{\hspace{3cm}}} \\
\multicolumn{2}{l}{\underline{\hspace{4cm}} \hfill \underline{\hspace{4cm}} \hfill \underline{\hspace{3cm}}} \\
\multicolumn{2}{l}{\underline{\hspace{4cm}} \hfill \underline{\hspace{4cm}} \hfill \underline{\hspace{3cm}}} \\
\multicolumn{2}{l}{\underline{\hspace{4cm}} \hfill \underline{\hspace{4cm}} \hfill \underline{\hspace{3cm}}} \\
\multicolumn{2}{l}{\underline{\hspace{4cm}} \hfill \underline{\hspace{4cm}} \hfill \underline{\hspace{3cm}}} \\
\multicolumn{2}{l}{(Attach additional sheet if more than 8 attendees)} \\
\midrule
\textbf{Presenter signature} & \underline{\hspace{6cm}} \\
\bottomrule
\end{tabularx}
\end{table}
